% TOP_PROBLEM: {"id": 10000088, "title": "Uniqueness at the percolation uniqueness threshold", "category_id": 19, "category": {"id": 19, "name": "probability", "display_name": "Probability", "description": "Problems involving probability theory, stochastic processes, and random structures.", "slug": "probability", "order_index": 19, "created_at": "2026-07-31T15:26:25.671Z"}, "status": "open", "problem_number": "AMR-099-0088", "set_id": 13, "statusQualification": "Retains only the endpoint classification requested in the catalogue; the original paper uses site percolation, and its update resolves the strictly-above-threshold part."}
% FIRST_POSTED: 2026-10-01
% ENTRY_KIND: statement_only
% UnsolvedMath ID 10000088; source catalogue code AMR-099-0088
% !TeX program = lualatex
% Definition and sources only; this file is not an LLM solution attempt.
\documentclass[11pt,a4paper]{article}
\usepackage[margin=25mm]{geometry}
\usepackage{fontspec}
\setmainfont{lmroman10-regular.otf}[BoldFont=lmroman10-bold.otf,
 ItalicFont=lmroman10-italic.otf,BoldItalicFont=lmroman10-bolditalic.otf]
\usepackage{amsmath,amssymb,mathtools}
\usepackage{unicode-math}
\setmathfont{latinmodern-math.otf}
\AtBeginDocument{\let\setminus\smallsetminus}
\usepackage{xurl,hyperref}
\hypersetup{colorlinks=true,urlcolor=blue}
\setcounter{secnumdepth}{0}
\setlength{\parindent}{0pt}
\setlength{\parskip}{5pt}
\setlength{\emergencystretch}{3em}
\begin{document}
\section*{Uniqueness at the percolation uniqueness threshold}\label{10000088}
{\small UnsolvedMath ID 10000088 · AMR-099-0088 · Probability · AMR Open Problem Lists}
\subsection{Definitions and mathematical statement}
Let $G$ be an infinite connected locally finite quasi-transitive graph, meaning that its automorphism group has finitely many orbits on vertices. In Bernoulli site percolation, retain each vertex independently with probability $p$ and consider the subgraph induced by the retained vertices. Denote its number of infinite connected components by $N_\infty(p)$, and put
\[
p_u(G)=\inf\{p\in[0,1]:\mathbb P_p(N_\infty(p)=1)=1\}.
\]
Characterize the graphs in this class for which
\[
\mathbb P_{p_u(G)}(N_\infty(p_u(G))=1)=1.
\]
The desired characterization should give necessary and sufficient structural conditions on the underlying graph, rather than restating the probability-one condition as a definition. The same classification question can be posed for Bernoulli bond percolation, with independent open edges.

This is specifically an endpoint question. The original list also asked for uniqueness for every $p>p_u$; its later update records that part as resolved. Uniqueness above the infimum does not determine whether it holds at the infimum itself. At the endpoint the alternative to exactly one infinite component can be absence of infinite components or the presence of infinitely many; a characterization must accommodate those distinct failures.

No planarity, amenability, Cayley-graph representation, or fixed degree is prescribed beyond local finiteness and quasi-transitivity. When $p_u=1$, site percolation retains the entire connected graph at the threshold, yielding a basic class with uniqueness. The difficult classification includes thresholds strictly below one and the effect of the graph's large-scale structure on critical uniqueness.
\subsection{Short English statement}
Which quasi-transitive graphs have a unique infinite percolation cluster exactly at the uniqueness threshold?
\subsection{Sources}
\begin{itemize}
\item[S1] ULAM.AI and the UnsolvedMath Contributors, supplied problems.json, record 10000088 (AMR-099-0088), AMR Open Problem Lists. Catalogue identity and source transcription; CC BY 4.0.
\url{https://huggingface.co/datasets/ulamai/UnsolvedMath}.
\item[S2] Benjamini and Schramm - Percolation beyond Z\textasciicircum{}d with 1999 update, Question 5, PDF page 10; at-threshold remainder, PDF page 10. Original source indexed by AMR-099-0088.
\url{https://emis.de/ft/41630#page=10}.
\item[S3] I. Benjamini and O. Schramm, Percolation beyond Z\textasciicircum{}d, many questions and a few answers, Questions 3 and 5; site-percolation convention in Section 2.
\url{https://emis.de/ft/41630}.
\item[S4] T. Hutchcroft and M. Pan, Percolation at the uniqueness threshold via subgroup relativization (2024), partial results for bond percolation on groups.
\url{https://arxiv.org/abs/2409.12283}.
\end{itemize}
\end{document}
